% AA-69
\chapter[Mianownik Weyla]{Elementy antyniezmiennicze i tożsamość mianownika Weyla}
\label{aa:ch:weyl-denominator}

W AA-68 wykazaliśmy, że charaktery nieprzywiedlne tworzą bazę pierścienia niezmienników grupy Weyla. Teraz zbadamy elementy zmieniające znak przy odbiciach prostych. Otrzymamy tożsamość mianownika, potrzebną do dowodu wzoru charakterów Weyla.

W całym rozdziale $K$ jest algebraicznie domkniętym ciałem charakterystyki zero, $\mathfrak g$ niezerową skończenie wymiarową półprostą algebrą Liego, $\mathfrak h$ jej maksymalną podalgebrą toralną, $\Phi$ układem korzeni, $\Delta=(\alpha_1,\ldots,\alpha_r)$ bazą, a $\Phi^+$ zbiorem korzeni dodatnich. Używamy kraty wag $P$, wag dominujących $P_+$ i porządku $\preceq$ z Definicji~\ref{aa:def:lie-dominant-integral}. Rachunki znaków wykonujemy w rzeczywistym modelu $E=P\otimes_{\mathbb Z}\mathbb R$ z AA-60, nie w ciele $K$. Przez $s_i$ oznaczamy odbicie odpowiadające $\alpha_i$, przez $W$ grupę Weyla, a przez $\ell$ funkcjonał $\ell(\alpha_i)=1$.

\section{Długość i orbity grupy Weyla}

\begin{lemat}
\label{aa:lem:weyl-positive-permutation}
Przy powyższych założeniach $s_i$ permutuje zbiór $\Phi^+\setminus\{\alpha_i\}$ i wysyła $\alpha_i$ na $-\alpha_i$. Jeżeli $w\in W$ przeprowadza wszystkie korzenie dodatnie na dodatnie, to $w=1$.
\end{lemat}
\begin{proof}
Jeżeli $\beta\in\Phi^+\setminus\{\alpha_i\}$, to co najmniej jeden współczynnik $\beta$ przy innym korzeniu prostym jest dodatni, gdyż układ jest zredukowany. Odbicie $s_i\beta=\beta-\beta(h_i)\alpha_i$ nie zmienia tego współczynnika. Otrzymany korzeń nie może być ujemny, więc jest dodatni. Nie może być równy $\alpha_i$, ponieważ jego obraz przez $s_i$ byłby ujemny. Z inwolutywności $s_i$ wynika żądana permutacja.

Udowodnimy pomocniczą zasadę skracania. Jeżeli $u=s_{i_1}\cdots s_{i_m}$ i $u\alpha_j$ jest korzeniem ujemnym, to, stosując kolejne odbicia od prawej, w pewnym kroku $k$ korzeń dodatni staje się ujemny. Niech $v=s_{i_{k+1}}\cdots s_{i_m}$. Z pierwszej części dowodu $v\alpha_j=\alpha_{i_k}$. Zgodnie ze wzorem na odbicie ortogonalne $v s_j v^{-1}=s_{i_k}$. Oznacza to $s_{i_k}v s_j=v$, a zatem
\[
 u s_j=s_{i_1}\cdots s_{i_{k-1}}v.
\]
Po prawej stronie pozostało $m-1$ odbić.

Niech teraz $w\ne1$ i wybierzmy najkrótszy zapis $w=u s_j$ w odbiciach prostych (istnieje na mocy Twierdzenia~\ref{aa:tw:weyl-chambers}). Gdyby $u\alpha_j$ był ujemny, zasada skracania dawałaby krótszy zapis $w$. Zatem $u\alpha_j$ jest dodatni, a $w\alpha_j=-u\alpha_j$ ujemny. Dowodzi to ostatniej tezy.
\end{proof}

\begin{lemat}
\label{aa:lem:weyl-regular-orbits}
Przy powyższych założeniach każda orbita $W$ w $P$ zawiera wagę dominującą. Jeżeli $\mu$ nie leży na żadnej hiperpłaszczyźnie korzeniowej w $E$, to orbita zawiera dokładnie jedną wagę dominującą, jest ona ściśle dominująca ($\mu(h_i)>0$ dla wszystkich $i$), a stabilizator każdego elementu tej orbity jest trywialny.
\end{lemat}
\begin{proof}
Skończoność $W$ pozwala wybrać w orbicie wagę $\nu$ o maksymalnej wartości $\ell$. Gdyby $\nu(h_i)<0$, to $\ell(s_i\nu)=\ell(\nu)-\nu(h_i)>\ell(\nu)$, sprzecznie z wyborem. Zatem $\nu\in P_+$. Jeżeli orbita jest regularna, czyli unika wszystkich hiperpłaszczyzn korzeniowych, to żadna z liczb $\nu(h_i)$ nie jest zerem.

W rzeczywistym modelu znak $\nu(h_i)$ jest znakiem $(\nu,\alpha_i)$. Jeżeli $\nu$ jest ściśle dominująca, to $(\nu,\beta)>0$ dla każdego $\beta\in\Phi^+$, gdyż $\beta$ jest niezerową nieujemną kombinacją korzeni prostych. Dla korzeni ujemnych znak jest przeciwny. Jeżeli $w\nu$ także jest ściśle dominująca, to dla $\beta\in\Phi^+$ mamy $(\nu,w^{-1}\beta)=(w\nu,\beta)>0$. Wobec tego $w^{-1}\beta\in\Phi^+$. Lemat~\ref{aa:lem:weyl-positive-permutation} daje $w^{-1}=1$. To dowodzi jednoznaczności i trywialności stabilizatora $\nu$. Stabilizatory pozostałych elementów orbity są z nim sprzężone, ponieważ $g(w\nu)=w\nu$ wtedy i tylko wtedy, gdy $w^{-1}gw\nu=\nu$.
\end{proof}

\section{Waga rho i elementy antyniezmiennicze}

\begin{lemat}
\label{aa:lem:weyl-rho}
Przy powyższych założeniach waga $\rho=\frac12\sum_{\alpha\in\Phi^+}\alpha$ jest sumą wag fundamentalnych $\omega_1+\cdots+\omega_r$. W szczególności $\rho\in P$ i jest ściśle dominująca. Waga $\nu\in P$ jest ściśle dominująca wtedy i tylko wtedy, gdy $\nu=\rho+\lambda$ dla $\lambda\in P_+$.
\end{lemat}
\begin{proof}
Z Lematu~\ref{aa:lem:weyl-positive-permutation} mamy $s_i\rho=\rho-\alpha_i$. Z definicji odbicia $s_i\rho=\rho-\rho(h_i)\alpha_i$. Ponieważ $\alpha_i\ne0$, dostajemy $\rho(h_i)=1$. Elementy $h_i$ są bazą $\mathfrak h$, więc $\rho=\sum_i\omega_i$. Dla wagi całkowitej warunek $\nu(h_i)>0$ oznacza $\nu(h_i)\ge1$. Jest to równoważne warunkowi $(\nu-\rho)(h_i)\ge0$ dla każdego $i$, czyli $\nu-\rho\in P_+$.
\end{proof}

\begin{definicja}
\label{aa:def:weyl-alternant}
Dla $w\in W$ oznaczmy $\epsilon(w)=\det_E w\in\{-1,1\}$. Element $F\in\mathbb Z[P]$ nazywamy antyniezmienniczym, jeśli $wF=\epsilon(w)F$ dla każdego $w\in W$. Dla $\nu\in P$ definiujemy alternant
\[
 A_\nu=\sum_{w\in W}\epsilon(w)e^{w\nu}.
\]
\end{definicja}

\begin{twierdzenie}
\label{aa:tw:weyl-alternant-basis}
Przy powyższych założeniach alternanty $A_{\rho+\lambda}$, $\lambda\in P_+$, tworzą bazę grupy abelowej elementów antyniezmienniczych pierścienia $\mathbb Z[P]$.
\end{twierdzenie}
\begin{proof}
Wyznacznik jest multiplikatywny, a $\epsilon(s_i)=-1$. Dla $g\in W$ podstawienie $u=gw$ w skończonej sumie daje $gA_\nu=\epsilon(g)A_\nu$. Dla elementu antyniezmienniczego $F=\sum_\mu c_\mu e^\mu$ porównanie współczynników daje $c_{w\mu}=\epsilon(w)c_\mu$. Jeżeli $\mu$ leży na hiperpłaszczyźnie korzeniowej $\alpha^\perp$, to odbicie $s_\alpha$ ustala $\mu$, a jego wyznacznik to $-1$. Mamy $c_\mu=-c_\mu$, zatem $c_\mu=0$ w $\mathbb Z$.

Nośnik $F$ jest więc sumą skończenie wielu orbit regularnych. Na mocy Lematu~\ref{aa:lem:weyl-regular-orbits} każda taka orbita zawiera dokładnie jedną wagę ściśle dominującą $\nu$ i każde $w\nu$ występuje dla dokładnie jednego $w$. Z Lematu~\ref{aa:lem:weyl-rho} mamy $\nu=\rho+\lambda$ dla pewnej $\lambda\in P_+$. Część $F$ wsparta na tej orbicie jest dokładnie $c_\nu A_\nu$, gdyż współczynnik przy $e^{w\nu}$ jest równy $\epsilon(w)c_\nu$. Sumowanie po orbitach daje rozpiętość. Różne alternanty mają rozłączne nośniki, a współczynnik przy $e^\nu$ w $A_\nu$ wynosi jeden, więc są liniowo niezależne nad $\mathbb Z$.
\end{proof}

\section{Tożsamość mianownika}

\begin{twierdzenie}[Tożsamość mianownika Weyla]
\label{aa:tw:weyl-denominator}
Przy powyższych założeniach w pierścieniu $\mathbb Z[P]$ zachodzi
\[
 e^\rho\prod_{\alpha\in\Phi^+}(1-e^{-\alpha})
 =\sum_{w\in W}\epsilon(w)e^{w\rho}=A_\rho.
\]
\end{twierdzenie}
\begin{proof}
Oznaczmy lewą stronę przez $D$. Z Lematu~\ref{aa:lem:weyl-rho} i permutacji dodatnich korzeni różnych od $\alpha_i$ otrzymujemy
\[
 \begin{aligned}
 s_iD&=e^{\rho-\alpha_i}(1-e^{\alpha_i})
       \prod_{\alpha\in\Phi^+\setminus\{\alpha_i\}}(1-e^{-\alpha})\\
 &=-e^\rho(1-e^{-\alpha_i})
       \prod_{\alpha\in\Phi^+\setminus\{\alpha_i\}}(1-e^{-\alpha})=-D.
 \end{aligned}
\]
Odbicia proste generują $W$, zatem $D$ jest antyniezmienniczy. Rozwinięcie iloczynu pokazuje, że każda waga nośnika ma postać $\rho-\sum_{\alpha\in S}\alpha$, gdzie $S\subseteq\Phi^+$. Zatem jest $\preceq\rho$, a współczynnik przy $e^\rho$ wynosi jeden, ponieważ niepusty zbiór dodatnich korzeni ma sumę o dodatniej wysokości.

Rozważmy dowolną orbitę występującą w nośniku $D$. Z dowodu Twierdzenia~\ref{aa:tw:weyl-alternant-basis} ma ona jedyną wagę ściśle dominującą $\nu=\rho+\lambda$, $\lambda\in P_+$. Ponieważ $\nu\preceq\rho$, mamy $\ell(\lambda)\le0$. Tymczasem w dowodzie Twierdzenia~\ref{aa:tw:lie-highest-classification} wykazano $|W|\lambda=\sum_i b_i\alpha_i$ z $b_i\in\mathbb Z_{\ge0}$. Jeżeli $\lambda\ne0$, to co najmniej jedno $b_i$ jest dodatnie, a więc $\ell(\lambda)=\sum_i b_i/|W|>0$. Zatem $\lambda=0$ i jedyną orbitą występującą w $D$ jest orbita $\rho$. Współczynnik przy $e^\rho$ jest jeden, więc $D=A_\rho$.
\end{proof}

\begin{przyklad}[Typy $A_1$ i $A_2$]
\label{aa:prz:weyl-denominator-small}
Dla typu $A_1$ mamy $\rho=\omega$, $\alpha=2\omega$, a po oznaczeniu $x=e^\omega$ mianownik jest równy $x(1-x^{-2})=x-x^{-1}$. Dla $A_2$ w bazie wag fundamentalnych mamy $\alpha_1=2\omega_1-\omega_2$, $\alpha_2=-\omega_1+2\omega_2$, $\rho=\omega_1+\omega_2$. Po oznaczeniu $x=e^{\omega_1}$, $y=e^{\omega_2}$ otrzymujemy
\[
 \begin{aligned}
 D&=xy(1-x^{-2}y)(1-xy^{-2})(1-x^{-1}y^{-1})\\
 &=xy-x^{-1}y^2-x^2y^{-1}+x^{-2}y+xy^{-2}-x^{-1}y^{-1}.
 \end{aligned}
\]
Dwa odbicia działają na współrzędnych przez $(p,q)\mapsto(-p,p+q)$ i $(p,q)\mapsto(p+q,-q)$. Orbita $(1,1)$ ma sześć elementów: $(1,1)$, $(-1,2)$, $(2,-1)$, $(-2,1)$, $(1,-2)$ i $(-1,-1)$. Znaki odpowiadają kolejno iloczynom złożonym z zera, jednego, jednego, dwóch, dwóch i trzech odbić. Otrzymujemy ten sam alternant.
\end{przyklad}

Mianownik $A_\rho$ jest niezerowy, gdyż współczynnik przy $e^\rho$ wynosi jeden. Następny rozdział wykaże, że dla $\lambda\in P_+$ zachodzi $(\operatorname{ch}L(\lambda))A_\rho=A_{\lambda+\rho}$. Sama antyniezmienniczość iloczynu nie wystarcza do tego wniosku: trzeba jeszcze wykluczyć pozostałe alternanty o niższych wagach.
